\documentclass[12pt]{ximera}
\title{Exam 2 (2023)}
\begin{document}
\begin{abstract}
An archived exam on Chapter 2 and the start of Chapter 3: gridlock, anarchy, and veto power, Banzhaf power, a mayor with a veto, counting sequential coalitions, continuous versus discrete assets, and divider-chooser.
\end{abstract}
\maketitle

\section*{Exam 2 (2023)}

This exam allowed one handwritten sheet of notes.

\begin{problem}
Consider the following weighted voting systems.
\[
\begin{array}{ll}
1) & $[35:9,8,7,6]$\\
2) & $[19:9,6,5,4]$\\
3) & $[13:11,9,6,2]$\\
4) & $[10:10,9,6,4]$\\
5) & $[18:7,4,2,1]$\\
\end{array}
\]

\begin{prompt}
\textbf{(a)} Which systems have gridlock?
\begin{selectAll}
\choice[correct]{System 1}
\choice{System 2}
\choice{System 3}
\choice{System 4}
\choice[correct]{System 5}
\end{selectAll}
\begin{feedback}
Gridlock means the quota exceeds the total weight, so no coalition can win.
System 1: $9+8+7+6=30<35$. System 5: $7+4+2+1=14<18$.
\end{feedback}
\end{prompt}

\begin{prompt}
\textbf{(b)} Which systems have anarchy?
\begin{selectAll}
\choice{System 1}
\choice{System 2}
\choice[correct]{System 3}
\choice[correct]{System 4}
\choice{System 5}
\end{selectAll}
\begin{feedback}
Anarchy means the quota is at most half the total weight, so two disjoint coalitions
could both win. System 3: $11+2\ge13$ and $9+6\ge13$. System 4: $10\ge10$ and $9+4\ge10$ (in fact $P_1$ is a dictator). System 2 is neither.
\end{feedback}
\end{prompt}

\begin{prompt}
\textbf{(c)} Which players in system 2 have veto power?
\begin{selectAll}
\choice[correct]{$P_1$}
\choice[correct]{$P_2$}
\choice{$P_3$}
\choice{$P_4$}
\end{selectAll}
\begin{feedback}
Without $P_1$ the others total $6+5+4=15<19$; without $P_2$, $9+5+4=18<19$. But $9+6+4=19$ and $9+6+5=20$, so $P_3$ and $P_4$ have no veto.
\end{feedback}
\end{prompt}
\end{problem}

\begin{problem}
Four business partners make decisions with the weighted voting system $[13:7,6,5,2]$, with
players $P_1,\dots,P_4$ in that order.

\begin{prompt}
\textbf{(a)} How many winning coalitions are there? $\answer{6}$
\begin{feedback}
$\set{P_1,P_2}=13$, $\set{P_1,P_2,P_3}=18$, $\set{P_1,P_2,P_4}=15$, $\set{P_1,P_3,P_4}=14$, $\set{P_2,P_3,P_4}=13$, and the grand coalition $20$.
\end{feedback}
\end{prompt}

\begin{prompt}
\textbf{(b)} Find the critical counts and the Banzhaf power distribution, as fractions.

Critical counts: $P_1$: $\answer{4}$, $P_2$: $\answer{4}$,
$P_3$: $\answer{2}$, $P_4$: $\answer{2}$

\smallskip
$\beta_1=\answer{1/3}$, $\beta_2=\answer{1/3}$,
$\beta_3=\answer{1/6}$, $\beta_4=\answer{1/6}$
\begin{feedback}
Both players are critical in $\set{P_1,P_2}$; $P_1,P_2$ in each of $\set{P_1,P_2,P_3}$ and $\set{P_1,P_2,P_4}$; all three in $\set{P_1,P_3,P_4}$ and in $\set{P_2,P_3,P_4}$; nobody in the grand coalition. The total is $12$: $\frac4{12},\frac4{12},\frac2{12},\frac2{12}$.
\end{feedback}
\end{prompt}
\end{problem}

\begin{problem}
A city council has $9$ seats: $8$ councilors and the mayor, $M$. They use simple majority ($5$ votes), but the mayor has a veto: no motion passes if the mayor votes No. Write
$\underline{\ }$ for a councilor.

\begin{prompt}
\textbf{(a)} In a sequential coalition with $M$ in slot $6$, which slot holds the
pivotal player? $\answer{6}$
\begin{feedback}
Five councilors cannot pass anything without the mayor, so the coalition first wins when $M$ joins in slot $6$.
\end{feedback}
\end{prompt}

\begin{prompt}
\textbf{(b)} In a sequential coalition with $M$ in slot $3$, which slot
holds the pivotal player? $\answer{5}$
\begin{feedback}
The first $5$ players include the mayor and form a majority; the councilor in slot $5$ is pivotal.
\end{feedback}
\end{prompt}

\begin{prompt}
\textbf{(c)} In which slots is the mayor pivotal?
\begin{selectAll}
\choice{1}
\choice{2}
\choice{3}
\choice{4}
\choice[correct]{5}
\choice[correct]{6}
\choice[correct]{7}
\choice[correct]{8}
\choice[correct]{9}
\end{selectAll}
\begin{feedback}
A coalition wins exactly when it has $5$ or more members including $M$. If $M$ is in slots $1$--$4$, the pivot is slot $5$; if $M$ is in slot $5$ or later, $M$ is the pivot. So the mayor is pivotal in slots $5$--$9$, a Shapley-Shubik power of $\frac59$.
\end{feedback}
\end{prompt}
\end{problem}

\begin{problem}
How many ways are there to fill slots $1,2,3,5,6,7,8$ (seven slots) of a sequential coalition
with seven different councilor names A, B, C, D, E, F, G? (This is unrelated to the
previous problem.) $\answer{5040}$
\begin{feedback}
Only the number of slots and names matters: $7!=7\cdot6\cdot5\cdot4\cdot3\cdot2\cdot1=5040$.
\end{feedback}
\end{problem}

\begin{problem}
Each of the following is to be divided among $5$ players. Is it a continuous or a
discrete fair-division problem?

\begin{prompt}
\textbf{(a)} A sandwich that is one part meat lover's, one part BLT, and one part veggie.
\begin{multipleChoice}
\choice[correct]{Continuous}
\choice{Discrete}
\end{multipleChoice}
\begin{feedback}
Continuous: a sandwich can be cut anywhere.
\end{feedback}
\end{prompt}
\begin{prompt}
\textbf{(b)} $75{,}000$ tons of uranium ore.
\begin{multipleChoice}
\choice[correct]{Continuous}
\choice{Discrete}
\end{multipleChoice}
\begin{feedback}
Continuous: any bulk quantity can be split in any proportion.
\end{feedback}
\end{prompt}
\begin{prompt}
\textbf{(c)} A sealed, signed original copy of a historic document.
\begin{multipleChoice}
\choice{Continuous}
\choice[correct]{Discrete}
\end{multipleChoice}
\begin{feedback}
Discrete: it cannot be divided without destroying its value.
\end{feedback}
\end{prompt}
\end{problem}

\begin{problem}
Jack and Jill have inherited a house and a plot of land. They both think the house is
worth \$160K. Jack thinks the land is worth \$180K, and Jill thinks it is
worth \$200K. The land can be divided in any way necessary. Jill is the divider and
makes share $s_1$ from the house and some of the land, and share $s_2$ from the rest of the
land.

\begin{prompt}
\textbf{(a)} What fraction of the land goes with the house in $s_1$, so that both shares
are equal in Jill's opinion? $\answer{1/10}$
\begin{feedback}
Jill values everything at $160+200=360$ (in thousands), so each share must be worth
$180$ to her. The house needs another $20$ in land, which is $\frac{20}{200}=\frac1{10}$
of the land. So $s_1=H+\frac1{10}L$ and $s_2=\frac9{10}L$.
\end{feedback}
\end{prompt}

\begin{prompt}
\textbf{(b)} How much is each share worth in Jack's opinion, in thousands of
dollars?

$s_1$: $\answer{178}$ \quad $s_2$: $\answer{162}$
\begin{feedback}
$s_1=160+\frac1{10}\cdot180=178$ and $s_2=\frac9{10}\cdot180=162$.
\end{feedback}
\end{prompt}

\begin{prompt}
\textbf{(c)} Which share will Jack choose?
\begin{multipleChoice}
\choice[correct]{$s_1$}
\choice{$s_2$}
\end{multipleChoice}
\begin{feedback}
$s_1$, worth $178>170$, half of Jack's total of $340$. Note that a house is indivisible: had Jill valued the house above half the total, she could not have made two equal shares this way.
\end{feedback}
\end{prompt}
\end{problem}

\end{document}
